\documentclass[11pt]{article}
\usepackage{amsmath,amssymb,amsthm}
\usepackage[margin=1in]{geometry}
\newtheorem{theorem}{Theorem}
\newtheorem{lemma}{Lemma}
\newtheorem{definition}{Definition}
\title{V102: Strong Affine Backdoors for Range Avoidance}
\author{NC0\_k-Avoid Laboratory}
\date{}
\begin{document}
\maketitle

\begin{definition}
Let $C:\{0,1\}^n\to\{0,1\}^m$ with $m>n$. A set $B$ of input variables is a
\emph{strong affine backdoor} if, for every assignment
$\sigma\in\{0,1\}^B$, every restricted output coordinate of $C|_\sigma$ is
affine over $\mathbb F_2$. Write $\beta(C)$ for the minimum size of such a set.
\end{definition}

\begin{theorem}[Strong-affine-backdoor avoider]
Given a strong affine backdoor $B$, one can construct deterministically a word
outside $\operatorname{Range}(C)$ in time
\[
  O\!\left(2^{|B|}\operatorname{poly}(|C|)\right).
\]
If every output has arity at most three, a backdoor of size at most $k$ can be
detected in $O(3^k\operatorname{poly}(|C|))$ time.
\end{theorem}

\begin{proof}
For an output prefix $p$, let
\[
 M(p)=|\{x:C(x)\text{ begins with }p\}|.
\]
Fix $\sigma$ on $B$. The prefix equations of $C|_\sigma$ form an affine linear
system, so Gaussian elimination gives its exact number of solutions. Summing
over all $2^{|B|}$ assignments gives $M(p)$ exactly. Starting from the empty
prefix, append the bit $b$ minimizing $M(pb)$. Since
$M(p0)+M(p1)=M(p)$, each step halves the count. Thus after $m$ bits,
\[
 M(y)\le 2^{n-m}<1,
\]
and integrality gives $M(y)=0$.

For detection at locality three, choose a gate that is not strongly affine
under the currently conditioned local coordinates. Strong-affine goodness is
upward closed. The inclusion-minimal good local supersets form an antichain in
a three-element support, hence there are at most three. Branch on those
supersets. Every branch adds at least one variable, giving a depth-at-most-$k$,
at-most-three-way search tree.
\end{proof}

\begin{lemma}[Exact residual-orbit rules]
For the $0x1b$ NPN orbit there is a unique selector coordinate; a local set is
strong-affine iff it contains the selector or both data coordinates. For the
signed-majority $0x17$ orbit, a local set is strong-affine iff it contains at
least two support coordinates.
\end{lemma}

\paragraph{Boundary.}
The theorem is parameterized by $\beta(C)$ and does not bound that parameter on
arbitrary circuits. It does not put unrestricted $NC^0_3$ range avoidance in
polynomial time, improve the unrestricted worst-case exponent, establish
novelty, or resolve P versus NP.
\end{document}
